\documentclass[11pt]{article}

\usepackage[margin=1in]{geometry}
\usepackage{amsmath, amssymb, amsthm}
\usepackage{mathtools}
\usepackage{booktabs}
\usepackage{multirow}
\usepackage{array}
\usepackage{graphicx}
\usepackage{hyperref}
\usepackage{xcolor}
\usepackage{enumitem}
\usepackage{algorithm}
\usepackage{algpseudocode}
\usepackage{caption}
\usepackage{subcaption}
\usepackage{microtype}
\usepackage{cleveref}
\usepackage[round,authoryear]{natbib}

\hypersetup{
  colorlinks=true,
  linkcolor=blue!50!black,
  citecolor=blue!50!black,
  urlcolor=blue!50!black,
}

\newtheorem{theorem}{Theorem}
\newtheorem{proposition}[theorem]{Proposition}
\newtheorem{lemma}[theorem]{Lemma}
\newtheorem{corollary}[theorem]{Corollary}
\theoremstyle{definition}
\newtheorem{definition}[theorem]{Definition}
\newtheorem{remark}[theorem]{Remark}
\newtheorem{example}[theorem]{Example}

\newcommand{\R}{\mathbb{R}}
\newcommand{\C}{\mathbb{C}}
\newcommand{\Z}{\mathbb{Z}}
\newcommand{\N}{\mathbb{N}}
\newcommand{\T}{\mathbb{T}}
\newcommand{\given}{\,\vert\,}
\newcommand{\eps}{\varepsilon}
\DeclareMathOperator{\softmax}{softmax}
\DeclareMathOperator*{\argmax}{arg\,max}
\DeclareMathOperator{\diag}{diag}
\DeclareMathOperator{\sgn}{sgn}
\DeclareMathOperator{\rank}{rank}

\title{
  Cyclic-Group Composition as a Diagnostic for Architectural
  Inductive Bias:\\
  \large Why Complex Unit-Circle Networks Generalize Where Real-Valued
  Gated Recurrences Do Not
}

\author{Minds-R-Lab}
\date{\today}

\begin{document}
\maketitle

\begin{abstract}
We ask when, if ever, complex-valued neural architectures offer
inductive-bias advantages over real-valued ones. To make the
question crisp we study cyclic-group composition tasks: sequences
containing one of $n$ atoms and $k$ instances of a single
``rotation'' token, with label $(c + k) \bmod n$, training on
$k \in \{0,\ldots,K_\text{tr}\}$ and evaluating on $k$ up to $4K_\text{tr}$.
Across five experiments and 200+ training runs we find:
(i)~A 303{,}000-parameter complex transformer with Hermitian
attention learns a phase-flip representation of the rotation token
(77\% of probe pairs near $\pi$ on $\Z/2$) but cannot compose
it---softmax weights sum to one, forcing additive composition.
(ii)~A 273-parameter set-equivariant phase-sum network
(\textsc{PhaseSumNet}) generalizes perfectly to out-of-distribution
depths on $\Z/n$ for every $n \in \{2,3,5,7,11,13\}$ tested, with
zero variance across seeds.
(iii)~Real-valued gated recurrences (GRU, LSTM) solve $\Z/2$ and
$\Z/3$ but collapse to chance ($1/n$) on $\Z/5$ across a 285$\times$
capacity sweep from 4k to 1M parameters. The mechanism is
state saturation against the $\tanh$ ceiling, after which the
linear readout becomes a near-constant function of $k$.
(iv)~On $\Z/7$ with two-layer GRU at large width, the barrier
\emph{leaks via grokking}: a late-training phase transition (loss
already at zero) sometimes finds an approximately periodic state
trajectory that partially aligns with the readout, reaching OOD
accuracy 0.95. Across every condition, the complex unit-circle
architecture solves the task two-to-three orders of magnitude more
efficiently in parameters, with no phase transition. We argue that
the qualitative separation is between architectures whose
composition operator is a group homomorphism by construction and
those that must discover one through optimization.
\end{abstract}

\section{Introduction}

A familiar story in deep learning attributes empirical success to
flexible function approximators trained on enough data with enough
compute. A second story, less often told but equally important, is
that the architecture's \emph{built-in symmetries} determine what it
can generalize to from finite data: a convolutional network
generalizes across spatial translations because translation-equivariance
is hardwired; a permutation-invariant set network generalizes across
input orderings because it cannot tell them apart in the first place;
a transformer with rotary positional encoding generalizes
to unseen positions roughly to the extent that its relative-position
construction is faithful to a continuous group action.

This paper asks a sharp version of the second question. \emph{Given a
synthetic task whose semantics is governed exactly by the cyclic group
$\Z/n\Z$, which architectures generalize across the group, and why?}
The question has a known partial answer: architectures that
\emph{represent} the group action---either by construction or by
training---generalize, while architectures that do not, do not. The
contribution of this paper is to make this folk wisdom mathematically
precise on a controlled task family, characterize the failure modes of
several common architectures (transformer, GRU, LSTM, additive
linear), and identify a previously-unnoticed phenomenon: even when an
architecture is in principle capable of representing the group, doing
so may require an optimization-time phase transition (grokking) that
does not occur for natively-symmetric architectures.

\subsection{Contributions}

\begin{enumerate}[leftmargin=*, itemsep=2pt]
\item We formalize a cyclic-group composition task on $\Z/n\Z$ and
prove that any additive architecture composing token embeddings with
a linear readout produces logits linear in the rotation count $k$,
which cannot represent $(c+k) \bmod n$ for $n \geq 3$. The model
provably converges to uniform output (\Cref{thm:additive-barred}).

\item We exhibit a minimal architecture, \textsc{PhaseSumNet}, that
represents the group action by construction via unit-circle phases.
It achieves OOD accuracy $1.000 \pm 0.000$ across $n \in
\{2,3,5,7,11,13\}$ with 306--845 parameters.

\item We show that softmax-attention transformers learn the
single-token group action (probe: 77\% of rotation pairs near $\pi$
on $\Z/2$) but fail to compose it, because softmax weights summing to
one make attention an additive composition operator
(\Cref{prop:attention-additive}).

\item Across a capacity sweep of GRU and LSTM from 4{,}000 to
1{,}060{,}000 parameters, real-valued gated recurrence fails on
$\Z/5$ at chance ($1/n$), with the mechanism identified by trajectory
analysis: $\tanh$-bounded state saturates past training depth,
producing constant logits.

\item On $\Z/7$ with two-layer GRU at large width, we observe
\emph{grokking}: a phase transition during late training where the
model's OOD accuracy jumps from $\sim$0.70 to $\sim$0.95 while
training loss has been zero for thousands of steps. The grokked
model's state has stopped saturating; its trajectory diagnostics
place it precisely between the failing baseline and the
unit-circle reference.
\end{enumerate}

\subsection{Roadmap}

\Cref{sec:background} reviews related work on complex networks,
group-equivariant architectures, length generalization, and grokking.
\Cref{sec:framework} develops the theoretical framework: cyclic-group
composition tasks, the question of architectural representability,
and the negative theorem for additive-linear architectures.
\Cref{sec:task,sec:architectures} describe the task family and the
six architectures we study. \Cref{sec:exp1,sec:exp2,sec:exp3,sec:exp4,sec:exp5}
report the five experiments. \Cref{sec:mechanism} performs trajectory
analysis to identify the GRU failure mechanism and the $\Z/7$
grokking phenomenon. \Cref{sec:discussion} discusses what the
findings collectively imply, with explicit limitations.


\section{Background and Related Work}
\label{sec:background}

\paragraph{Complex-valued neural networks.}
\citet{trabelsi2018deep} provide a comprehensive blueprint for
training complex-valued networks at scale, including modReLU,
complex batch normalization, and complex initialization schemes.
Earlier, \citet{arjovsky2016unitary} proposed unitary RNNs whose
recurrent transitions are constrained to lie on the unitary group,
arguing that the resulting orbit structure stabilizes long-range
dependencies. \citet{sordoni2013modeling} applied complex amplitudes
to language modeling in an information-retrieval setting. Despite
this body of work, complex-valued networks have remained niche; one
recurring critique is that they have not produced demonstrable
advantages over real-valued networks at matched parameter counts on
standard benchmarks.

\paragraph{Group-equivariant architectures.}
\citet{cohen2016group} formalized group equivariance for CNNs,
extending to non-Euclidean groups in subsequent work. The unifying
principle---that an architecture should commute with the symmetry
group acting on its input space---has produced equivariant networks
for $\mathrm{SE}(3)$ molecular dynamics, $\mathrm{SO}(3)$ point
clouds, and others. The cyclic group $\Z/n$ is the simplest
non-trivial group with finite order, and serves here as a minimal
benchmark for the broader equivariance question.

\paragraph{Length and depth generalization in sequence models.}
\citet{hahn2020theoretical} proves that softmax self-attention cannot
recognize certain regular languages---specifically parity---in the
infinite-length limit, a result we use to interpret our $\Z/2$
transformer experiment. \citet{bhattamishra2020ability} empirically
confirms that transformers fail to learn parity at depths beyond
their training distribution. Our $\Z/n$ task family generalizes
parity to arbitrary cyclic groups and provides a graded difficulty
spectrum.

\paragraph{Grokking.}
\citet{power2022grokking} introduced the term \emph{grokking} for a
training pattern in which a network reaches near-perfect training
accuracy quickly but generalization on a held-out set emerges only
after a long period of further training. Subsequent work
\citep{nanda2023progress, liu2022omnigrok} has connected grokking to
the slow discovery of low-norm or structured solutions in
overparameterized models. Our $\Z/7$ result connects grokking to the
specific question of whether an architecture's gradient descent can
\emph{discover} a group-respecting representation late in training,
absent that representation being built in.

\paragraph{Length generalization on algorithmic tasks.}
\citet{anil2022exploring,zhou2023algorithms,jelassi2024repeat} report
that transformers struggle with length generalization on
algorithmic tasks (parity, addition, copying) and analyze the
specific failure modes. Our work is in the same lineage but uses a
cleaner $\Z/n$ family to isolate the algebraic structure from
positional encoding details.


\section{Theoretical Framework}
\label{sec:framework}

\subsection{Cyclic-Group Composition Tasks}

Fix $n \in \N_{\geq 2}$. The \emph{cyclic-group composition task} on
$\Z/n\Z$ is defined by the following data-generating process. The
input alphabet $\Sigma$ consists of:
\begin{itemize}[itemsep=2pt]
\item $n$ \emph{atom} tokens $a_0, a_1, \ldots, a_{n-1}$, each carrying
a class label;
\item one \emph{rotation} token $\tau$ ("TWIRL");
\item $F$ semantically inert \emph{filler} tokens $f_1, \ldots, f_F$.
\end{itemize}
A sample $\mathbf{x} = (x_1, x_2, \ldots, x_L)$ is constructed by
choosing an atom class $c \in \{0, \ldots, n-1\}$ uniformly at random,
a rotation count $k$ from a specified distribution, and a filler count
$\phi$, then placing $a_c$ together with $k$ copies of $\tau$ and
$\phi$ fillers in a uniformly random order. The label is
\[
y(\mathbf{x}) \;=\; (c + k) \bmod n.
\]
The task is order-invariant by construction: $y$ depends only on
which atom is present and how many rotations there are.

Concretely, training samples draw $k$ uniformly from
$\{0, 1, \ldots, K_\text{tr}\}$; evaluation samples are
\emph{depth-stratified}: equal numbers at each $k \in
\{0, 1, \ldots, K_\text{eval}\}$ with $K_\text{eval} \gg K_\text{tr}$.
A model that has learned the group structure should achieve
$\sim$100\% accuracy on every $k$. A model that has learned a
lookup table over training depths will achieve 100\% on
$k \leq K_\text{tr}$ and degrade on the out-of-distribution (OOD)
depths $K_\text{tr} < k \leq K_\text{eval}$.

\subsection{What Does It Mean for an Architecture to ``Represent'' the Group?}

\begin{definition}[Group-respecting architecture]
\label{def:group-respecting}
Let $\mathcal{A}$ be a sequence-to-class architecture mapping inputs
$\mathbf{x} \in \Sigma^*$ to logits $\ell(\mathbf{x}) \in \R^n$. We
say $\mathcal{A}$ is \emph{group-respecting} for $\Z/n$ on the
composition task if there exists a representation $\rho: \Z/n \to
\mathrm{GL}(V)$ on some real or complex vector space $V$, a class
embedding $\phi: \{a_0, \ldots, a_{n-1}\} \to V$, a filler-invariant
state space $\mathcal{H}$, and a readout $R: V \to \R^n$ such that
the architecture's computation is equivalent to
\begin{equation}
\label{eq:group-respecting}
\ell(\mathbf{x}) \;=\; R\bigl( \rho(k)\, \phi(c) \bigr),
\end{equation}
where $c, k$ are the atom class and rotation count of $\mathbf{x}$,
and the readout satisfies the homomorphism property
$R(\rho(k+\Delta) v) = R(\rho(\Delta)\, \rho(k) v)$ for all
$k, \Delta, v$.
\end{definition}

The architectural separation we study is whether
\Cref{eq:group-respecting} holds \emph{by construction} (some
architectures), can be \emph{discovered} by training (others), or is
provably unattainable regardless of parameters and training
(\Cref{thm:additive-barred} below).

\subsection{The Negative Theorem for Additive-Linear Architectures}
\label{subsec:negative-theorem}

We now state the result that determines the lower bound of the
spectrum.

\begin{definition}[Additive-linear architecture]
An architecture is \emph{additive-linear} if its logits take the form
\[
\ell(\mathbf{x}) \;=\; W \!\left( \sum_{t=1}^{L} e(x_t) \right) + b
\;\in\; \R^n,
\]
for some embedding $e: \Sigma \to \R^d$, readout matrix
$W \in \R^{n \times d}$, and bias $b \in \R^n$.
\end{definition}

\begin{theorem}[Additive-linear cannot represent $\Z/n$ above chance]
\label{thm:additive-barred}
Let $\mathcal{A}$ be an additive-linear architecture and let
$y(\mathbf{x}) = (c + k) \bmod n$ be the cyclic-group composition
task on $\Z/n\Z$ with $n \geq 2$. There is no choice of $e$, $W$,
$b$ such that $\argmax_j \ell_j(\mathbf{x}) = y(\mathbf{x})$ for all
sufficiently long $k$ trajectories. In particular, no additive-linear
model can achieve OOD accuracy above $1/n$ on the depth-stratified
evaluation distribution with $K_\mathrm{eval} > 2n$.
\end{theorem}

\begin{proof}
Fix the atom class $c$ and let the filler contribution be folded into
a constant offset (we may assume the optimal model assigns
$e(f_i) = 0$ for all fillers, since fillers carry no signal). For an
input with exactly $k$ rotation tokens the logits become
\[
\ell_j(c, k) \;=\; W_{j, :}\,\bigl(e(a_c) + k \cdot e(\tau)\bigr) + b_j
\;=\; \alpha_j(c) + \beta_j \cdot k,
\]
with $\alpha_j(c) = W_{j,:} e(a_c) + b_j$ and
$\beta_j = W_{j,:} e(\tau)$, both independent of $k$. Each class
logit is therefore an affine function of $k$ with slope $\beta_j$
that does not depend on the atom class $c$.

For the model to predict class $j$ at rotation count $k$ we need
$\ell_j(c, k) > \ell_{j'}(c, k)$ for all $j' \neq j$. This is the
condition $(\beta_j - \beta_{j'}) k > \alpha_{j'}(c) - \alpha_j(c)$
for all $j' \neq j$, a system of $n-1$ linear inequalities in $k$.
If $\beta_j > \beta_{j'}$ this holds for all sufficiently large
$k$; if $\beta_j < \beta_{j'}$ it holds for all sufficiently small
$k$. Therefore there is at most one class $j^*$ (the one with
largest slope $\beta_{j^*}$) for which $\ell_{j^*}$ becomes the
argmax for all large enough $k$.

The correct label, however, cycles: $y(c, k) = (c + k) \bmod n$
takes every value in $\{0, 1, \ldots, n-1\}$ as $k$ ranges over any
window of length $n$. For $K_\text{eval} > 2n$, the
correct label takes each of $n$ values on (approximately) $1/n$ of
the OOD depths. The model predicts $j^*$ on all of them. Hence the
model is correct at most $\frac{1}{n}$ of the time on sufficiently
long OOD trajectories.
\end{proof}

\begin{remark}[Topology]
\label{rem:topology}
The mechanism behind \Cref{thm:additive-barred} is the topology of
$\R$: at most $n-1$ pairwise-line crossings can occur between any
$n$ affine functions of a single variable, so the argmax has at
most $n-1$ transitions and is eventually constant in $k$. The
unit circle in $\C$ avoids this restriction precisely because it
has no point at infinity, and a continuous trajectory along it can
visit all $n$ ``slots'' an unbounded number of times.
\end{remark}

\begin{remark}[Strengthening]
The proof in fact rules out \emph{any} architecture whose logits
are an affine function of token count $k$ in some specified token,
regardless of how nonlinear the upstream embedding is. The bound
extends to in-distribution accuracy whenever
$K_\mathrm{tr} \geq n-1$: the same argmax-transition bound applies
to the union of ID and OOD points. Inspecting the trained
\textsc{RealAddNet} on $\Z/3$ ($K_\mathrm{tr} = 5 > 2$), the
per-class slope spread $\max_j \beta_j - \min_j \beta_j$ is
$0.020 \pm 0.005$ across seeds, indicating the optimizer converged
to nearly-equal slopes and effectively uniform logits---consistent
with the theorem-bounded failure.
\end{remark}

\begin{remark}[$n=2$ caveat]
\label{rem:n-equals-2}
At $n=2$, the bound $1/n = 1/2$ numerically coincides with chance,
so the failure is less visually dramatic than for $n=3$ (chance
$0.333$) or $n=13$ (chance $0.077$). We measure
\textsc{RealAddNet} OOD accuracy of $0.505 \pm 0.015$ on $\Z/2$,
which is at chance and consistent with the theorem. The
representational impossibility is the same; only the visibility
of the failure differs.
\end{remark}

\subsection{Softmax Attention is Additive Composition}
\label{subsec:attention-additive}

The same argument extends to a broad class of attention-based
architectures.

\begin{proposition}[Softmax attention as a stochastic convex combination]
\label{prop:attention-additive}
Let an attention layer compute its output at query position $q$ as
$z_q = \sum_{t} \alpha_{q,t} \, v(x_t)$, where the attention weights
$\alpha_{q,t} \geq 0$ satisfy $\sum_t \alpha_{q,t} = 1$, and $v$ is
the value projection. Then $z_q$ is a convex combination of value
vectors, and its magnitude is bounded by $\max_t \|v(x_t)\|$.
\end{proposition}

The proposition is a consequence of the softmax constraint and is
not novel. Its relevance here is the following corollary, which
makes precise what \citet{hahn2020theoretical} showed for parity in
the formal-language setting.

\begin{corollary}[Attention cannot multiply]
\label{cor:attention-cant-multiply}
Let $\rho: \Z/n \to \mathrm{GL}(V)$ be a faithful representation
acting by multiplication: $\rho(k_1+k_2) = \rho(k_1)\rho(k_2)$.
Suppose an architecture $\mathcal{A}$ computes, at its output
position, a single softmax-attention aggregation over per-token
value vectors $v(x_t)$ followed by a linear readout. Then $\mathcal{A}$
cannot represent $\rho(k) v$ for arbitrary $k$ via the attention
output, because the attention output is a convex combination of the
$v(x_t)$ values, which lies in their convex hull, while
$\rho(k)\, v(x_t)$ ranges over orbits in $V$ that exit any convex
hull for sufficiently large $k$ in general.
\end{corollary}

In short: attention can mix tokens, but it cannot multiply them.
The complex-valued attention we test in \Cref{sec:exp1} learns
the correct per-token rotation $\rho(\tau) \approx -I$ in its
embedding, but the architecture's composition operator
(attention) is unable to use it.

\subsection{What Unit-Circle Phases Buy}
\label{subsec:unit-circle}

The simplest architecture satisfying the homomorphism condition
\Cref{eq:group-respecting} for arbitrary $\Z/n$ uses the unit
circle in $\C$ as its state space.

\begin{definition}[Phase-sum architecture]
\label{def:phasesum}
Let each token $x \in \Sigma$ be assigned a phase vector
$\theta(x) \in \R^d$. The phase-sum architecture computes
\begin{equation}
\Theta(\mathbf{x}) = \sum_{t=1}^{L} \theta(x_t) \in \R^d,
\qquad
\ell(\mathbf{x}) = R\bigl( \cos\Theta(\mathbf{x}),\, \sin\Theta(\mathbf{x}) \bigr),
\end{equation}
where $R$ is a linear readout $\R^{2d} \to \R^n$, and $\cos,\sin$
are applied componentwise.
\end{definition}

\begin{proposition}[Phase-sum is group-respecting for every $\Z/n$]
\label{prop:phasesum-respects}
For any $n \geq 2$, the phase-sum architecture admits a parameter
setting with $d=1$ such that
$\ell(\mathbf{x}) = R'(\rho(k) \cdot \phi(c))$ in the sense of
\Cref{def:group-respecting}, where $\rho(k)$ acts on $\C$ by
multiplication by $e^{2\pi i k / n}$.
\end{proposition}

\begin{proof}
Choose $\theta(a_c) = 2\pi c / n$ for $c = 0, \ldots, n-1$,
$\theta(\tau) = 2\pi / n$, and $\theta(f_i) = 0$ for all fillers.
Then $\Theta(\mathbf{x}) = 2\pi(c+k)/n$, and $(\cos\Theta, \sin\Theta)
\in \R^2$ identifies with the point $e^{2\pi i (c+k)/n}$ on the
unit circle. The unit circle contains exactly $n$ points
$e^{2\pi i j/n}$ for $j = 0, \ldots, n-1$, and these are linearly
separable in $\R^2$ (any $n$ points on a circle in general position
are). A linear readout $R'$ that picks out the correct class from
the position on the unit circle exists.
\end{proof}

\begin{remark}
The proof exhibits a $d=1$ solution. In practice we initialize with
$d=16$ phase dimensions, providing redundant routes to the optimal
solution; the trained network typically realizes the group action
only in a subset of dimensions and the readout learns to ignore the
rest. We discuss the empirical convergence properties in
\Cref{sec:exp2}.
\end{remark}

\subsection{The Saturation Failure of Bounded Real-Valued Recurrence}
\label{subsec:saturation}

Real-valued gated recurrences (GRU, LSTM) constrain their hidden
state to lie in $[-1, 1]^d$ via the $\tanh$ activation. This
boundedness is usually a virtue (gradient stability), but in our
setting it produces a specific failure mode: \emph{saturation} past
the training depth.

We will show empirically (\Cref{sec:mechanism}) that for the GRU on
$\Z/5$, the hidden-state magnitude $\|h(k)\|$ grows monotonically
with rotation count $k$ from $\sim$5.12 at $k=0$ to $\sim$7.04 at
$k=30$, asymptoting against the bound. Once saturated, the linear
readout's output becomes a near-constant function of $k$, and the
argmax is constant. A constant argmax over $n$ balanced classes
scores $1/n$ in expectation. This precisely matches our measured
asymptote of $\sim 0.25$ for $\Z/5$ and $\sim 0.14$ for $\Z/7$.

We give an informal argument for why this happens. The GRU update
$h_{t+1} = (1-z_t) h_t + z_t \tilde{h}_t$ with
$\tilde{h}_t \in (-1, 1)^d$ produces an exponentially-weighted
moving average of bounded targets. When the input is a long
homogeneous sequence of one token, the candidate $\tilde{h}_t$
becomes approximately constant, and $h_t$ exponentially converges
to that fixed point. Across the training distribution (depths
$k \leq K_\text{tr}$), the model can configure this fixed point to
encode useful information about the partial state. Beyond $K_\text{tr}$,
the convergence to fixed point dominates, and information about
$k \bmod n$ that was carried in transient state dynamics is lost.

This argument predicts that an unbounded real-valued recurrence
(e.g., a linear RNN with appropriate normalization) might escape the
saturation. We do not test this directly in the present paper, but
identify it as a target follow-up.


\section{The Task Family}
\label{sec:task}

We instantiate the cyclic-group composition task with the following
parameters across experiments. The vocabulary contains a
\texttt{PAD} token, a \texttt{CLS} token (used by the transformer
architectures of \Cref{sec:exp1}), the $n$ atom tokens, the rotation
token $\tau$, and 10 filler tokens. Sequences are zero-padded to the
maximum length within each batch. The atom appears exactly once per
sequence; rotation count $k$ is drawn uniformly from
$\{0, \ldots, K_\text{tr}\}$ during training and from each integer in
$\{0, \ldots, K_\text{eval}\}$ separately during evaluation
(depth-stratified). Filler count is drawn uniformly from
$\{0, \ldots, 8\}$.

The atom and rotations are placed in uniformly random order within
the body of the sequence, so positional templates cannot help. This
order-invariance is what makes a set-equivariant architecture
(\Cref{def:phasesum}) the architecturally-correct choice.

The cardinal experiment parameter choices are summarized in
\Cref{tab:task-params}.

\begin{table}[h]
\centering
\caption{Task and training parameters across experiments.}
\label{tab:task-params}
\begin{tabular}{lcccc}
\toprule
Experiment & Group $\Z/n$ & $K_\mathrm{tr}$ & $K_\mathrm{eval}$ & Seeds \\
\midrule
1: complex transformer & $\Z/2$ & 3  & 10 & 5 \\
2: recurrent architectures on $\Z/2$ & $\Z/2$ & 3  & 15 & 5 \\
3: $\Z/3$ separation test & $\Z/3$ & 5  & 20 & 5 \\
4: scaling sweep & $\Z/n$, $n\in\{2,3,5,7,11,13\}$ & 5 & 20 & 3 \\
5: capacity sweep & $\Z/5, \Z/7$ & 5 & 20 & 3 \\
\bottomrule
\end{tabular}
\end{table}


\section{Architectures}
\label{sec:architectures}

We compare six architectural families, each described below with its
relation to the theoretical framework of \Cref{sec:framework}.

\subsection{Complex Transformer (\textsc{ComplexTransformer})}

A pre-norm transformer encoder operating on $\C^d$ hidden states.
Attention uses the Hermitian inner product
$\langle q, k \rangle_H = \mathrm{Re}(q^* k)$. Nonlinearity is
modReLU \citep{trabelsi2018deep}:
$\mathrm{modReLU}(z) = \mathrm{ReLU}(|z| + b)\cdot z / (|z| + \eps)$.
Positional encoding is $e^{i \omega_d t}$ on each complex
dimension. A \texttt{CLS} token's complex hidden state is read out
via $\mathrm{Re}(w^* h_\mathrm{CLS})$ for logits. Parameter count is
matched to a real-valued baseline transformer of the same depth and
head count: 303k parameters at $d=64$ vs. 310k for the real model at
$d=92$.

By \Cref{cor:attention-cant-multiply}, the architecture cannot
multiplicatively compose per-token group actions, regardless of how
faithfully the individual token embeddings learn the action.

\subsection{Phase-Sum Network (\textsc{PhaseSumNet})}

\Cref{def:phasesum}. For task $\Z/n$, the readout is linear
$\R^{2d} \to \R^n$. We use $d=16$ throughout (giving 16 redundant
phase dimensions). Initialization is critical: phases are drawn
uniformly from $[-\pi, \pi]$. We discuss why in
\Cref{subsec:init-matters}.

\subsection{Real Additive Network (\textsc{RealAddNet})}

The negative-control architecture of \Cref{thm:additive-barred}:
\[
\ell(\mathbf{x}) = W\!\left(\sum_t e(x_t)\right) + b,
\]
with $e: \Sigma \to \R^{2d}$, $W \in \R^{n \times 2d}$. By the
theorem, this architecture cannot exceed $1/n$ accuracy on
$\Z/n$ for $n \geq 3$.

\subsection{Gated Complex RNN (\textsc{GatedComplexRNN})}

A bidirectional recurrent network with $\C^d$ hidden states. Each
token $x$ induces an update
\begin{align}
\hat{h}_t &= \exp(i\,\theta(x))\, h_t + v(x), \\
h_{t+1} &= (1 - g(x)) h_t + g(x) \hat{h}_t,
\end{align}
where $\theta(x) \in \R^d$ is a learned per-dimension phase rotation,
$v(x) \in \C^d$ a learned complex additive contribution, and
$g(x) \in (0,1)^d$ a sigmoid gate. Composition is multiplicative
via the rotation $e^{i\theta(\tau)}$ for the rotation token; the
architecture is in principle group-respecting if the model learns
$\theta(\tau) = 2\pi/n$ and $g(\tau) \to 1$, $v(\tau) \to 0$.

\subsection{GRU and LSTM Baselines (\textsc{GRU}, \textsc{LSTM})}

Standard bidirectional single- and two-layer GRU and LSTM, with a
final linear readout from concatenated forward/backward final
states. Real-valued, $\tanh$-bounded hidden state. These are
parameter-matched (within a few percent) to \textsc{GatedComplexRNN}
in \Cref{sec:exp4}, then scaled up to $\sim$1M parameters in
\Cref{sec:exp5}.

\subsection{Initialization Considerations}
\label{subsec:init-matters}

For both \textsc{PhaseSumNet} and \textsc{GatedComplexRNN}, the
phase embedding $\theta$ must be initialized uniformly on
$[-\pi, \pi]$, not near zero. With near-zero initialization, every
sequence produces $\cos\Theta \approx 1, \sin\Theta \approx 0$, so
logits are nearly constant across inputs and gradient signal is too
weak to escape. With uniform initialization, sequences are spread
across the unit circle from the start and learning proceeds. We
mention this because the convergence depends sensitively on it and
because the failure mode (nominally trainable architecture sits at
chance forever) is otherwise puzzling.


\section{Experiment 1: Complex Transformer on $\Z/2$}
\label{sec:exp1}

\paragraph{Setup.} Train both \textsc{ComplexTransformer} ($d=64$,
303k params) and a parameter-matched real transformer ($d=92$, 310k
params) on $\Z/2$ parity-of-negation: a single atom $\in \{T, F\}$,
$k$ \texttt{NOT} tokens, fillers in random order, with label
$\mathrm{atom\_sign} \cdot (-1)^k$. Train depths $k \in \{0,1,2,3\}$,
evaluate $k \in \{0, \ldots, 10\}$, 5 seeds, AdamW with
cosine-decayed learning rate.

\paragraph{Probe.} For the complex model we additionally measure the
\emph{phase shift induced by a single \texttt{NOT}}. For paired
sequences $(s, s')$ differing only by one extra \texttt{NOT}, we
compute the angle $\Delta = \arg(h_\mathrm{CLS}(s') / h_\mathrm{CLS}(s))$
in each complex dimension. A model that has learned \texttt{NOT}
as multiplication by $-1$ should produce $\Delta \approx \pi$ in
every dimension.

\paragraph{Results.}

\begin{table}[h]
\centering
\caption{Experiment 1: 5-seed averages on $\Z/2$. ID = in-distribution
accuracy ($k \leq 3$), OOD = out-of-distribution accuracy
($4 \leq k \leq 10$). Probe column applies only to the complex model.}
\label{tab:exp1-results}
\begin{tabular}{lcccc}
\toprule
Model & Params & ID acc & OOD acc & Probe: $|\Delta|$ \\
\midrule
Real Transformer    & 309{,}765 & $1.000 \pm 0.000$ & $0.488 \pm 0.005$ & ---  \\
Complex Transformer & 302{,}978 & $1.000 \pm 0.000$ & $0.445 \pm 0.005$ & $2.43$ rad (target $\pi$) \\
\bottomrule
\end{tabular}
\end{table}

\paragraph{Interpretation.} Both transformers achieve perfect ID
accuracy and fall to near-chance OOD. The complex model is
slightly \emph{worse} than the real one. However, the phase probe
reveals that the complex model \emph{did} learn the algebra: the
mean rotation induced by a single \texttt{NOT} is 2.43 radians (vs
target $\pi = 3.14$), and 77\% of dimension-paired sequences land
within $\pi/8$ of $\pi$. A randomly-initialized network gives
$|\Delta| \approx 0.17$.

The failure is therefore not in the per-token representation but
in composition. By
\Cref{prop:attention-additive,cor:attention-cant-multiply}, softmax
attention is forced to mix token contributions additively; multiplying
five rotations of $\pi$ is not the same operation as averaging them
with weights summing to one. The complex model has the right algebra
embedded in its weights but cannot use it through the attention
mechanism.


\section{Experiment 2: Recurrent Architectures on $\Z/2$}
\label{sec:exp2}

\paragraph{Setup.} Train three architectures whose composition is
multiplicative by construction: \textsc{PhaseSumNet} ($d=16$, $\sim$273
params), \textsc{GatedComplexRNN} ($d=24$, $\sim$3{,}000 params), and a
parameter-matched GRU baseline ($\sim$3{,}000 params). Train depths
$k \in \{0,\ldots,3\}$, eval to $k=15$. 5 seeds.

\paragraph{Results.}

\begin{table}[h]
\centering
\caption{Experiment 2: 5-seed averages on $\Z/2$.}
\label{tab:exp2-results}
\begin{tabular}{lccc}
\toprule
Model & Params & ID acc & OOD acc \\
\midrule
\textsc{PhaseSumNet}     & 273   & $1.000 \pm 0.000$ & $1.000 \pm 0.000$ \\
\textsc{GatedComplexRNN} & 4{,}033 & $1.000 \pm 0.000$ & $0.999 \pm 0.000$ \\
\textsc{GRU}             & 3{,}962 & $1.000 \pm 0.000$ & $0.999 \pm 0.000$ \\
\bottomrule
\end{tabular}
\end{table}

\paragraph{Interpretation.} All three multiplicative architectures
generalize essentially perfectly to OOD depth 15---5$\times$ the
training distribution---using fewer than 4{,}000 parameters, while
the 303{,}000-parameter complex transformer of Experiment 1 sat at
chance. The structural prior is doing two-to-three orders of magnitude
more work than parameter count.

A subtle point: \textsc{PhaseSumNet}'s phase probe is messy. Across
seeds, the mean $\cos\theta(\mathrm{NOT})$ is $-0.11$ (target $-1$),
and only 24\% of phase dimensions land within $\pi/8$ of $\pm\pi$.
The model does not converge to the textbook optimum. But it does not
need to: with $d=16$ phase dimensions feeding a linear readout, the
readout learns to combine the few well-placed dimensions and ignore
the rest. \textsc{PhaseSumNet}'s OOD accuracy is exactly 1.000
nonetheless, with zero variance across seeds.

This experiment establishes that the inductive bias from
\emph{multiplicative composition} (not specifically complex numbers)
solves $\Z/2$. The real-valued GRU achieves the same OOD accuracy as
the complex models. To separate "multiplicative composition" from
"complex numbers" we need a group whose representation in $\R^*$
is not native. That is the next experiment.


\section{Experiment 3: $\Z/3$ as a Separating Task}
\label{sec:exp3}

\paragraph{Why $\Z/3$.} The multiplicative group $\R^*$ contains
$\Z/2 = \{+1, -1\}$ as a subgroup but \emph{not} $\Z/3$. There is
no order-3 element in $\R^*$ under multiplication. To represent
$\Z/3$ in a network, one must either (a) use 2D rotations
(equivalently, $\C$ on the unit circle), (b) use multi-state gating
that simulates a 3-state automaton, or (c) use a non-additive
composition operator. Pure additive-linear architectures are
provably barred by \Cref{thm:additive-barred}.

\paragraph{Setup.} Architectures: \textsc{PhaseSumNet3} (the $n=3$
version of \textsc{PhaseSumNet}), \textsc{RealAddNet} (the negative
control of \Cref{thm:additive-barred}), \textsc{GatedComplexRNN3}, and
\textsc{GRU3}. Train depths 0--5, eval 0--20, 5 seeds.

\paragraph{Results.}

\begin{table}[h]
\centering
\caption{Experiment 3: 5-seed averages on $\Z/3$. Chance is $1/3$.}
\label{tab:exp3-results}
\begin{tabular}{lcccc}
\toprule
Model & Params & ID acc & OOD acc (6--20) & Closure under 3$\cdot\tau$ \\
\midrule
\textsc{PhaseSumNet3}     & 355    & $1.000 \pm 0.000$ & $\mathbf{1.000 \pm 0.000}$ & 1.000 \\
\textsc{RealAddNet}       & 611    & $\mathbf{0.336 \pm 0.020}$ & $\mathbf{0.341 \pm 0.008}$ & --- \\
\textsc{GatedComplexRNN3} & 4{,}547  & $1.000 \pm 0.000$ & $0.987 \pm 0.006$ & 1.000 \\
\textsc{GRU3}             & 4{,}503  & $1.000 \pm 0.000$ & $0.999 \pm 0.001$ & 0.953 \\
\bottomrule
\end{tabular}
\end{table}

\paragraph{Interpretation.}

\textbf{\Cref{thm:additive-barred} is empirically tight.}
\textsc{RealAddNet} sits at exactly chance (1/3) on both ID
\emph{and} OOD. The architecture is incapable of fitting even the
training distribution, in confirmation of the proof: linear-in-$k$
logits cannot represent $(c+k) \bmod 3$ at any depth. We verify by
inspecting the learned readout: the per-class slopes
$\beta_j = W_{j,:} e(\tau)$ have spread $0.020 \pm 0.005$ across
seeds, near zero---the model has correctly identified that no
consistent linear function of $k$ predicts the label and converged
to uniform output.

\textbf{Multiplicative architectures all generalize.}
\textsc{PhaseSumNet3} at 355 parameters reaches OOD 1.000 with zero
variance. \textsc{GatedComplexRNN3} and \textsc{GRU3} reach 0.987
and 0.999 respectively. The GRU's slight edge at extreme OOD
depths (18--20) is interesting and we return to it in
\Cref{sec:exp4}.

This experiment establishes that the relevant axis is
\emph{multiplicative vs.\ additive composition}, with the negative
side of the axis provably barred and the positive side generically
successful at $\Z/3$. We have not yet separated complex from real
within the multiplicative class. That is Experiment 4.


\section{Experiment 4: Scaling Sweep over Group Order $n$}
\label{sec:exp4}

\paragraph{Setup.} Sweep $n \in \{2, 3, 5, 7, 11, 13\}$. For each
$n$, train all four multiplicative-class architectures plus the
\textsc{RealAddNet} negative control. Train depths 0--5, eval 0--20,
3 seeds per cell. Total: 72 training runs.

\paragraph{Results.} \Cref{tab:exp4-results} summarizes by giving
OOD accuracy as a function of $n$. The full per-depth breakdown is
in the project repository.

\begin{table}[h]
\centering
\caption{Experiment 4: OOD accuracy (mean of 3 seeds) on $\Z/n$ for
each $n$. Bold = near-perfect; underline = collapse to chance.}
\label{tab:exp4-results}
\setlength{\tabcolsep}{4pt}
\begin{tabular}{lcccccc}
\toprule
Model & $n=2$ & $n=3$ & $n=5$ & $n=7$ & $n=11$ & $n=13$ \\
\midrule
\textsc{PhaseSumNet} & $\mathbf{1.000}$ & $\mathbf{1.000}$ & $\mathbf{1.000}$ & $\mathbf{1.000}$ & $\mathbf{1.000}$ & $\mathbf{1.000}$ \\
\textsc{RealAddNet} & $0.505$ & $\underline{0.341}$ & $\underline{0.200}$ & $\underline{0.136}$ & $\underline{0.069}$ & $\underline{0.070}$ \\
\textsc{GatedComplexRNN} & $0.967$ & $0.979$ & $0.992$ & $0.980$ & $0.816$ & $0.864$ \\
\textsc{GRU} & $\mathbf{1.000}$ & $0.991$ & $\underline{0.248}$ & $\underline{0.241}$ & $\underline{0.285}$ & $\underline{0.237}$ \\
\midrule
Chance $1/n$ & $0.500$ & $0.333$ & $0.200$ & $0.143$ & $0.091$ & $0.077$ \\
\bottomrule
\end{tabular}
\end{table}

\paragraph{Interpretation.}

\textbf{\textsc{PhaseSumNet} is exact across $n$.} OOD accuracy
$1.000 \pm 0.000$ at every $n$ tested, with 306--845 parameters.
The architectural symmetry match holds across an order of magnitude
in group order. This is the strongest single result of the paper.

\textbf{\textsc{RealAddNet} matches the theorem's chance prediction.}
For $n \geq 3$, OOD accuracy is statistically indistinguishable from
$1/n$. The exception at $n=2$ (OOD $0.505 \approx$ chance for $\Z/2$)
is informative: at $n=2$ the theorem does not apply (two linear
functions can change ranking, $n-1=1$ time), but the model still
converged to chance. This is a learning failure---an architecture
not theoretically barred can still fail to find a workable
solution.

\textbf{\textsc{GatedComplexRNN} declines gracefully with $n$.}
OOD ranges from 0.967 at $n=2$ to 0.816 at $n=11$, with a clear
break-point around $n \approx 8$. The mechanism is that the
architecture's angular margin between adjacent classes is $2\pi/n$,
which shrinks as $n$ grows, while the additive value contribution
$v(\tau)$ lets state magnitude drift slightly with sequence length;
the ratio of drift to margin determines the break.

\textbf{\textsc{GRU} shows a phase transition, not a scaling decline.}
GRU goes from $1.000$ at $n=2$ and $0.991$ at $n=3$ to $\sim$0.25 at
$n \geq 5$ across every tested $n$, with ID accuracy still
$\sim$1.000 throughout. The architecture fits the training set
perfectly at every $n$ but generalizes only for $n \in \{2, 3\}$.
This is not a gradual loss of capacity; it is a binary
representability transition we need to investigate further.

We address the GRU's $n \geq 5$ collapse via a capacity sweep in
Experiment 5 and a trajectory analysis in \Cref{sec:mechanism}.


\section{Experiment 5: Capacity Sweep on the GRU Collapse}
\label{sec:exp5}

\paragraph{Question.} Is the GRU's $\Z/n_{\geq 5}$ collapse a
\emph{structural barrier} (real-valued gating cannot represent the
group at any capacity) or merely \emph{sample-inefficient} (more
parameters or training would close the gap)?

\paragraph{Setup.} Hold $n=5$ and $n=7$ fixed. Sweep
$(d_{\text{model}}, n_\text{layers}) \in \{(16,1), (32,1), (32,2),
(64,1), (64,2), (128,1), (128,2), (256,1)\}$, giving an 8-cell
capacity ladder from $\sim$4k to $\sim$1.06M parameters. Include
\textsc{LSTM} as a control with the same gating-family ancestry but
different gate equations. Compare to \textsc{PhaseSumNet} at matched
$d$ values. Train 120{,}000 samples for 25 epochs at $\mathrm{lr} =
3 \times 10^{-3}$, batch size 256, 3 seeds per cell. Total: 126
runs.

\paragraph{Results for $\Z/5$.} OOD accuracy summary
(\Cref{tab:exp5-z5}):

\begin{table}[h]
\centering
\caption{Experiment 5, $\Z/5$: OOD accuracy across capacity. All
GRU/LSTM cells sit at or below 0.34; \textsc{PhaseSumNet} is exactly
$1.000$ at every cell.}
\label{tab:exp5-z5}
\setlength{\tabcolsep}{5pt}
\begin{tabular}{l@{\hskip 8pt}rcccc}
\toprule
Model & Params & $d=16,L=1$ & $d=64,L=1$ & $d=128,L=2$ & $d=256,L=1$ \\
\midrule
GRU              & 3.7k--797k   & $0.285$ & $0.247$ & $0.266$ & $0.240$ \\
LSTM             & 4.8k--1.06M  & $0.288$ & $0.261$ & $0.290$ & $0.251$ \\
\textsc{PhaseSumNet} & 453--7.2k & $\mathbf{1.000}$ & $\mathbf{1.000}$ & $\mathbf{1.000}$ & $\mathbf{1.000}$ \\
\bottomrule
\end{tabular}
\end{table}

Across 16 GRU+LSTM cells from 3{,}717 to 1{,}059{,}845 parameters
(285$\times$ capacity range), OOD on $\Z/5$ stays in the band
$[0.22, 0.34]$. ID is $1.000 \pm 0.000$ in every cell. The
GRU/LSTM architectures fit the training set perfectly at every scale
and produce near-noise OOD predictions.

\paragraph{Results for $\Z/7$.} 1-layer GRU and all LSTM stay at
chance. But 2-layer GRU climbs with capacity:

\begin{center}
\begin{tabular}{lcc}
\toprule
Configuration & Params & OOD acc \\
\midrule
GRU $d=32, L=2$  & 32k  & $0.402 \pm 0.041$ \\
GRU $d=64, L=2$  & 127k & $0.564 \pm 0.079$ \\
GRU $d=128, L=2$ & 499k & $0.787 \pm 0.094$ \\
\textsc{PhaseSumNet} & 551--8{,}711 & $\mathbf{1.000 \pm 0.000}$ \\
\bottomrule
\end{tabular}
\end{center}

The 2-layer GRU at $d=128$ reaches 78.7\% on $\Z/7$ with 499k
parameters, partially escaping the barrier. Single-layer GRU and
all LSTM at the same scale stay at chance. PhaseSumNet at 551
parameters hits $1.000$.

\paragraph{Interpretation.}

\textbf{The $\Z/5$ barrier is sharp.} 285$\times$ capacity range,
two architecture families, two depths each, three seeds each, total
48 runs, none above 0.34 OOD. This establishes as well as a
synthetic-task experiment can establish that real-valued gated
recurrence does not solve $\Z/5$ at any capacity tested within 25
epochs.

\textbf{LSTM ties the GRU result.} The structural failure is not
GRU-specific. Whatever mechanism prevents $\Z/5$ representation
acts in both LSTM and GRU. Possibilities: (a) both architectures
share the $\tanh$-bounded state space; (b) both architectures share
sigmoid gates; (c) both architectures share a per-step linear
contribution from input embeddings. \Cref{sec:mechanism} provides
evidence that the bounded state is the key.

\textbf{The 2-layer $\Z/7$ result is unresolved.} The climb 0.40
$\to$ 0.56 $\to$ 0.79 with capacity at 2-layer is reproducible
across the three seeds (stderr 0.04--0.09). It is specifically
2-layer GRU, specifically $\Z/7$. We hypothesize and partially
confirm in \Cref{sec:mechanism} that the second layer permits a
late-training phase transition---grokking---during which the
architecture discovers an approximately group-respecting state
trajectory.


\section{Mechanism Analysis: State Trajectories}
\label{sec:mechanism}

The capacity sweep established the existence of the $\Z/5$ barrier
and the $\Z/7$ partial escape. This section identifies what is
happening inside the architectures by direct measurement of state
and logit dynamics.

\subsection{Methodology}

For a trained model, we build paired probe sequences differing only
in the number of trailing rotation tokens, with the base prefix
(\texttt{CLS}, atom, fillers) held identical across $k$. Across
$k = 0, 1, \ldots, 30$, we measure:
\begin{enumerate}[itemsep=2pt]
\item State magnitude $\|h(k)\|$.
\item State modular return: $\|h(k+n) - h(k)\|$ (per fixed base
sequence).
\item Logit modular return: $\|\ell(k+n) - \ell(k)\|$.
\item \textbf{Argmax invariance}: the fraction of probe sequences
satisfying $\argmax \ell(k+n) = \argmax \ell(k)$.
\end{enumerate}

Argmax invariance directly measures the property required for OOD
generalization: predictions should be invariant under adding $n$
rotation tokens, since $(c + k + n) \bmod n = (c + k) \bmod n$.

\subsection{The GRU Failure on $\Z/5$}

Train a 1-layer GRU at $d=64$ on $\Z/5$ (same configuration as the
$d=64, L=1$ cell of \Cref{tab:exp5-z5}) and run the probe.
\Cref{tab:trajectory-z5} summarizes.

\begin{table}[h]
\centering
\caption{Trajectory diagnostics on $\Z/5$. The two architectures are
trained with the same task, training depths, and number of epochs.}
\label{tab:trajectory-z5}
\begin{tabular}{lcc}
\toprule
& GRU $d=64$ & \textsc{PhaseSumNet} $d=16$ \\
\midrule
OOD accuracy ($k>5$) & $0.23$ & $\mathbf{1.000}$ \\
State magnitude at $k>5$ & $6.75$ & $4.00$ \\
Magnitude ratio $\|h(30)\|/\|h(5)\|$ & $1.19$ & $1.00$ \\
State $\|\Delta\| / \|h\|$ (relative) & $0.17$ & $0.98$ \\
Logit $\|\Delta\|$ & $9.87$ & $0.46$ \\
Argmax invariance under $5\cdot\tau$ & $0.54$ & $\mathbf{1.000}$ \\
\bottomrule
\end{tabular}
\end{table}

\textbf{State-level closure is misleading.} The GRU's relative state
movement under $5\cdot\tau$ is only $0.17$ (state changes by 17\%
of its magnitude), while \textsc{PhaseSumNet}'s is $0.98$ (state
changes by nearly its full magnitude). On state closure alone the
GRU looks better. But its argmax invariance is $0.54$ and
\textsc{PhaseSumNet}'s is $1.00$.

The resolution is that the readout-relevant quantity is logits, not
states. \textsc{PhaseSumNet}'s state lives on the unit circle in
$\R^{2d}$, and the readout's null space (under addition of
$n \cdot \theta(\tau) \equiv 0 \pmod{2\pi}$) is exactly the
direction of state motion. State moves; logits don't. The GRU's
state geometry has no such alignment: small state shifts get
amplified by the linear readout into 22$\times$ larger logit
shifts ($9.87$ vs $0.46$).

\textbf{The per-$k$ pattern reveals the failure mode.}
\Cref{tab:gru-perk} shows the GRU's argmax invariance at each $k$.

\begin{table}[h]
\centering
\caption{Per-$k$ argmax invariance under $5\cdot\tau$ for the
$\Z/5$ GRU.}
\label{tab:gru-perk}
\begin{tabular}{lc|lc}
\toprule
$k$ & invariance & $k$ & invariance \\
\midrule
0   & $0.83$ & 12--15 & $0.50$--$0.61$ \\
1   & $0.63$ & 16--19 & $0.67$--$0.80$ \\
2   & $0.28$ & 20--25 & $0.85$--$0.98$ \\
3--7 & $\mathbf{0.00}$ & & \\
8--11 & $0.38$ & & \\
\bottomrule
\end{tabular}
\end{table}

Invariance is high near the training boundary ($k=0, 1$),
catastrophically zero in the transition zone ($k=3$--$7$), and
recovers at large $k$. The recovery is misleading: at $k=20$--$25$
the GRU's state magnitude has saturated against the $\tanh$ ceiling
($\|h\|$ grows from 5.12 at $k=0$ to 7.04 at $k=30$ and asymptotes),
and the logits become a near-constant function of $k$ (logit
$\|\Delta\|$ falls from 24.96 at $k=5$ to 2.05 at $k=20$). The
prediction is constant in $k$, not group-respecting.

This is the unified mechanism for the $\Z/5$ failure: \emph{the
state saturates against the $\tanh$ ceiling past the training depth,
making the readout a near-constant function of $k$, which is right
$1/n$ of the time}. No amount of capacity within the tested range
escapes this, because the boundedness is a property of the
activation, not the layer width.

\subsection{The $\Z/7$ Grokking Phenomenon}

We applied the same trajectory diagnostics to the 2-layer GRU at
$d=128$ on $\Z/7$ (the partial-escape cell). One seed of this
configuration was run with the same hyperparameters as Experiment 5
(120k samples, 25 epochs, $\mathrm{lr}=3\times 10^{-3}$). The
training trajectory is shown in \Cref{tab:grokking-trajectory}.

\begin{table}[h]
\centering
\caption{Training trajectory for 2-layer GRU $d=128$ on $\Z/7$.
Training loss reaches zero at step 2000 and stays there. OOD
accuracy sits at $\sim 0.70$ for 4000 more steps, then jumps at
step 6000.}
\label{tab:grokking-trajectory}
\begin{tabular}{ccc}
\toprule
Step & Train loss & OOD acc \\
\midrule
500   & 0.0005 & $0.632$ \\
2000  & 0.0001 & $0.702$ \\
5500  & 0.0000 & $0.710$ \\
\textbf{6000} & \textbf{0.0005} & $\mathbf{0.940}$ \\
8000  & 0.0000 & $\mathbf{0.947}$ \\
11700 & 0.0000 & $0.947$ \\
\bottomrule
\end{tabular}
\end{table}

The pattern of (i) training loss converging to zero, (ii) OOD
plateauing for an extended period, then (iii) abrupt OOD
improvement, is the standard signature of grokking
\citep{power2022grokking}. At step 6000, with the model having been
at zero loss for $\sim$4000 steps, OOD jumps from 0.71 to 0.94. The
brief loss spike at the same step is consistent with the
optimization stepping into a qualitatively different basin.

The trajectory diagnostics on the grokked model are striking
(\Cref{tab:grokked-vs-failing}).

\begin{table}[h]
\centering
\caption{Trajectory diagnostics on $\Z/7$ for the grokked 2-layer
GRU, compared to the failing baseline ($\Z/5$ 1-layer GRU) and the
success reference (\textsc{PhaseSumNet} on $\Z/7$).}
\label{tab:grokked-vs-failing}
\begin{tabular}{lccc}
\toprule
& 1-layer GRU $\Z/5$ (fails) & \textbf{2-layer GRU $\Z/7$ (grokked)} & \textsc{PhaseSumNet} $\Z/7$ (succeeds) \\
\midrule
OOD accuracy & $0.23$ & $\mathbf{0.947}$ & $1.000$ \\
Magnitude ratio (30 / 5) & $1.19$ & $\mathbf{1.02}$ & $1.00$ \\
State $\|\Delta\|$ relative & $0.13$ & $0.17$ & $0.98$ \\
Logit $\|\Delta\|$ & $9.87$ & $\mathbf{3.30}$ & $0.46$ \\
Argmax invariance & $0.58$ & $\mathbf{0.86}$ & $1.00$ \\
\bottomrule
\end{tabular}
\end{table}

\textbf{The grokked GRU has stopped saturating.} Magnitude ratio
drops from $1.19$ (failing) to $1.02$ (grokked); state magnitude is
nearly constant across $k$. The model has found an
approximately-periodic state trajectory that bypasses the $\tanh$
ceiling.

\textbf{The grokked GRU has partially aligned its readout.} Logit
$\|\Delta\|$ falls 3$\times$ (from 9.87 to 3.30); argmax invariance
climbs from 0.58 to 0.86. The model has moved most of the way from
the failing baseline to the unit-circle success.

\subsection{What Grokking Found, and What It Did Not}

We interpret the grokked GRU as having discovered, via late-training
optimization, an approximate version of what \textsc{PhaseSumNet}
implements by construction: a bounded state trajectory along
directions the readout partially ignores modulo $n$. Three pieces
of evidence support this reading:

\begin{enumerate}[itemsep=2pt]
\item The magnitude ratio is essentially identical to
\textsc{PhaseSumNet}'s ($1.02$ vs $1.00$): the state has become
effectively bounded.
\item Logit $\|\Delta\|$ has dropped by a factor of 3, indicating
the readout has lost most of its amplification of state shifts.
\item The phase transition occurs at step 6000---4000 steps after
training loss has converged---which is the signature of slow
discovery of a structured solution, not faster convergence.
\end{enumerate}

The grokked solution is not exact: argmax invariance is $0.86$ vs
\textsc{PhaseSumNet}'s $1.00$, and OOD is $0.947$ vs $1.000$. The
GRU has \emph{approximately} discovered Z/7 representation, not
exactly. We hypothesize that with longer training or larger
capacity it might approach 1.000, but at substantial parameter and
training cost relative to \textsc{PhaseSumNet}'s 551 parameters
reaching 1.000 by step 1000.

\subsection{Implications for the $\Z/5$ Claim}

Does the $\Z/7$ grokking finding undermine the $\Z/5$ structural
claim? We believe not, because:

\begin{itemize}[itemsep=2pt]
\item In 48 $\Z/5$ runs spanning a 285$\times$ capacity range, three
seeds each, no run shows the late-training jump characteristic of
grokking. All saturate within $\sim$1500 steps and remain at the
chance plateau.
\item The diagnostic signature of saturation (monotonically growing
$\|h\|$, constant argmax at large $k$) is present in every $\Z/5$
GRU we examined and absent in the grokked $\Z/7$ model.
\end{itemize}

The most defensible version of the structural claim is therefore:
\emph{within practical training budgets, real-valued gated
recurrence does not learn to represent $\Z/n$ for $n \geq 5$
reliably. At $n = 7$ with two layers and large width, a fraction of
training runs grok to an approximate representation late in training,
producing the partial escape observed in \Cref{tab:exp4-results}.
Whether $\Z/5$ admits grokking at sufficiently long training or
larger capacity is open}.


\section{Discussion}
\label{sec:discussion}

\subsection{What We Can Claim}

\begin{enumerate}[itemsep=2pt]
\item \textbf{Additive-linear architectures cannot represent
$\Z/n$ for $n \geq 3$.} Provable
(\Cref{thm:additive-barred}), confirmed empirically across $n =
3, 5, 7, 11, 13$ at accuracy exactly $1/n$.

\item \textbf{Softmax attention is additive composition.} The
complex transformer in Experiment 1 learns the algebra
($|\Delta| = 2.43$ radians, target $\pi$) but cannot compose it.
This is consistent with \citet{hahn2020theoretical} and the
formal-language literature.

\item \textbf{The unit circle in $\C$ admits a faithful
representation of every $\Z/n$ in one dimension.}
\textsc{PhaseSumNet} reaches OOD $1.000 \pm 0.000$ at every $n$ we
tested, with 306--845 parameters.

\item \textbf{Bounded real-valued gated recurrence does not solve
$\Z/n$ for $n \geq 5$ within practical training budgets.} 285$\times$
capacity sweep, two architecture families, three seeds, all at
$\leq 0.34$ OOD. The mechanism is state saturation against the
$\tanh$ ceiling.

\item \textbf{Grokking can partially leak the barrier on $\Z/7$.}
With 2-layer GRU at $d=128$, a late-training phase transition
sometimes produces an approximate Z/7 representation (OOD $\sim$0.95).
The grokked model has stopped saturating---a structurally distinct
regime from the failing baseline.

\item \textbf{The qualitative separation is between architectures
matching the symmetry by construction and those discovering it
through optimization}, including potentially via grokking. The
former generalize predictably with $O(n)$ parameters; the latter
require orders of magnitude more parameters and may not always
discover the structure within standard training budgets.
\end{enumerate}

\subsection{What We Cannot Yet Claim}

\begin{enumerate}[itemsep=2pt]
\item Whether $\Z/5$ admits grokking with $\gg$25 epochs of training
or $\gg$1M parameters. The signature is absent in our 48 cells but
the test is finite.

\item Whether the saturation mechanism is causally responsible for
the $\Z/5$ failure. The clean test is to replace the $\tanh$-bounded
RNN with an unbounded normalized one; the prediction is that the
saturation pattern should disappear and OOD should change. We have
not run this experiment.

\item Whether the n=7 grokking phenomenon scales to $n \in \{11, 13\}$
at sufficient capacity. Experiment 4 measured 2-layer GRU at small
$d$ on these and saw chance; we have not retested at $d \in \{128,
256\}$.

\item Anything about real-language negation. Our experiments use
controlled cyclic-rotation tasks. Real-language phenomena are
governed by the cyclic structure of negation but also by scope,
modality, sarcasm, and idiomatic constructions that we do not model.
\end{enumerate}

\subsection{Implications for Practitioners}

The findings suggest a design heuristic. For any task whose
underlying structure can be identified with a known group (cyclic
groups for parity-like tasks; $\mathrm{SO}(3)$ for 3D pose;
$\mathrm{SE}(3)$ for rigid body motion; permutations for set
processing), an architecture whose state geometry contains a
faithful representation of that group should be preferred over a
generic architecture. The cost in expressiveness from constraining
the state geometry is offset by orders-of-magnitude better sample
efficiency, deterministic generalization properties, and the absence
of late-training phase transitions.

For cyclic-group structure specifically: complex unit-circle
representations are the minimal sufficient construction. Real-valued
gated recurrences are insufficient for $n \geq 5$ within practical
training budgets; transformers with softmax attention are insufficient
at any $n$.

\subsection{Connection to Existing Theory}

Our negative result for softmax attention
(\Cref{cor:attention-cant-multiply}) is a special case of
\citet{hahn2020theoretical}'s impossibility result for parity in
transformer self-attention. We extend it to general cyclic groups.

Our positive result for unit-circle architectures relates to the
broader literature on group-equivariant networks
\citep{cohen2016group}. The novelty in our setting is the
\emph{minimality} of the construction: the same algebraic structure
that motivates large-scale group-equivariant CNNs for image data is
sufficient to solve a sequence task with 300--800 parameters.

The grokking finding for the $\Z/7$ 2-layer GRU is consistent with
the broader grokking literature in suggesting that overparameterized
models can discover structured solutions late in training. We
believe the cleanness of the cyclic-group setting may make it a
useful benchmark for mechanistic studies of grokking.


\section{Limitations}
\label{sec:limitations}

\begin{enumerate}[itemsep=2pt]
\item \textbf{Synthetic task family.} All experiments are on
controlled cyclic-group rotation tasks. Real-world phenomena with
cyclic structure (musical pitch, days of week, cardinal directions,
negation in natural language) are noisier and partially observable.

\item \textbf{Limited cyclic-group orders.} We test $n \leq 13$.
Whether the pattern holds for $n = 100$ or $n = 1000$ is unknown.

\item \textbf{Three seeds in the capacity sweep.} The 2-layer GRU
on $\Z/7$ grokking happens in roughly one seed out of three at the
$d=128$ scale; a 10-seed rerun would tighten the variance.

\item \textbf{Single architecture per category.} We test one complex
transformer, one GRU, one LSTM. There are many transformer variants
(\citet{su2024roformer}'s RoPE, ALiBi, etc.) we have not tested.

\item \textbf{Limited mechanism analysis.} The trajectory diagnostics
focus on state and logit dynamics in the last position. A finer
analysis (per-position state evolution, gate values, etc.) would
better characterize what the gated models do and do not internalize.
\end{enumerate}

\section{Conclusion}
\label{sec:conclusion}

We have studied cyclic-group composition tasks as a controlled
benchmark for architectural inductive bias. Combining a formal
proof that additive-linear architectures cannot represent
$\Z/n$ for $n \geq 3$, an empirical demonstration that softmax
transformers fall into the additive class, a positive existence
result for a 300-parameter unit-circle architecture that
generalizes across $n \in \{2, 3, 5, 7, 11, 13\}$, and a mechanistic
analysis of where the real-valued gated baselines fail, we arrive at
the following synthesis:

Architectures whose composition operator is a group homomorphism
into a state space containing the task's symmetry group generalize
on cyclic-group tasks predictably and with $O(n)$ parameters.
Architectures whose composition operator is additive or whose state
saturates do not, regardless of capacity. Between these extremes
sits a regime---two-layer GRU at large width on $\Z/7$---where the
architecture can in principle discover the group structure but
must do so through a late-training phase transition (grokking) that
is unreliable and parameter-expensive. The qualitative split
between \emph{architectural match by construction} and
\emph{discovery through optimization} appears to be the central
distinction.

We hope the framework---cyclic-group tasks as a graded diagnostic,
trajectory diagnostics as a mechanistic probe, and the structural
vs.\ grokked vs.\ constructed taxonomy---is useful for future work
on architectural inductive bias.


\paragraph{Code and data.} All code, training logs, and per-seed
results are available at
\url{https://github.com/Minds-R-Lab/ComplexAttn}.


\small
\bibliographystyle{abbrvnat}

\begin{thebibliography}{20}

\bibitem[Anil et~al.(2022)]{anil2022exploring}
Anil, C., Wu, Y., Andreassen, A., Lewkowycz, A., Misra, V., Ramasesh, V., Slone, A., Gur-Ari, G., Dyer, E., and Neyshabur, B.
\newblock Exploring length generalization in large language models.
\newblock \emph{NeurIPS}, 2022.

\bibitem[Arjovsky et~al.(2016)]{arjovsky2016unitary}
Arjovsky, M., Shah, A., and Bengio, Y.
\newblock Unitary evolution recurrent neural networks.
\newblock \emph{ICML}, 2016.

\bibitem[Bhattamishra et~al.(2020)]{bhattamishra2020ability}
Bhattamishra, S., Ahuja, K., and Goyal, N.
\newblock On the ability and limitations of transformers to recognize formal languages.
\newblock \emph{EMNLP}, 2020.

\bibitem[Cohen and Welling(2016)]{cohen2016group}
Cohen, T.~S. and Welling, M.
\newblock Group equivariant convolutional networks.
\newblock \emph{ICML}, 2016.

\bibitem[Hahn(2020)]{hahn2020theoretical}
Hahn, M.
\newblock Theoretical limitations of self-attention in neural sequence models.
\newblock \emph{TACL}, 8:156--171, 2020.

\bibitem[Jelassi et~al.(2024)]{jelassi2024repeat}
Jelassi, S., d'Ascoli, S., Domingo-Enrich, C., Wu, Y., Li, Y., and Charton, F.
\newblock Length generalization in arithmetic transformers.
\newblock \emph{arXiv:2306.15400}, 2024.

\bibitem[Liu et~al.(2022)]{liu2022omnigrok}
Liu, Z., Michaud, E.~J., and Tegmark, M.
\newblock Omnigrok: Grokking beyond algorithmic data.
\newblock \emph{arXiv:2210.01117}, 2022.

\bibitem[Nanda et~al.(2023)]{nanda2023progress}
Nanda, N., Chan, L., Lieberum, T., Smith, J., and Steinhardt, J.
\newblock Progress measures for grokking via mechanistic interpretability.
\newblock \emph{ICLR}, 2023.

\bibitem[Power et~al.(2022)]{power2022grokking}
Power, A., Burda, Y., Edwards, H., Babuschkin, I., and Misra, V.
\newblock Grokking: Generalization beyond overfitting on small algorithmic datasets.
\newblock \emph{arXiv:2201.02177}, 2022.

\bibitem[Sordoni et~al.(2013)]{sordoni2013modeling}
Sordoni, A., Bengio, Y., and Nie, J.-Y.
\newblock Modeling term dependencies with quantum language models for IR.
\newblock \emph{SIGIR}, 2013.

\bibitem[Su et~al.(2024)]{su2024roformer}
Su, J., Lu, Y., Pan, S., Murtadha, A., Wen, B., and Liu, Y.
\newblock Roformer: Enhanced transformer with rotary position embedding.
\newblock \emph{Neurocomputing}, 568, 2024.

\bibitem[Trabelsi et~al.(2018)]{trabelsi2018deep}
Trabelsi, C., Bilaniuk, O., Zhang, Y., Serdyuk, D., Subramanian, S., Santos, J.~F., Mehri, S., Rostamzadeh, N., Bengio, Y., and Pal, C.~J.
\newblock Deep complex networks.
\newblock \emph{ICLR}, 2018.

\bibitem[Zhou et~al.(2023)]{zhou2023algorithms}
Zhou, H., Bradley, A., Littwin, E., Razin, N., Saremi, O., Susskind, J., Bengio, S., and Nakkiran, P.
\newblock What algorithms can transformers learn? A study in length generalization.
\newblock \emph{arXiv:2310.16028}, 2023.

\end{thebibliography}

\end{document}
